\documentclass[11pt]{article}
\usepackage[margin=1in]{geometry}
\usepackage{amsmath,amssymb,amsthm}
\usepackage{graphicx}
\usepackage[hidelinks]{hyperref}
\setlength{\emergencystretch}{2em}

\newtheorem{theorem}{Theorem}
\newtheorem{corollary}{Corollary}
\newtheorem{definition}{Definition}
\newtheorem{remark}{Remark}

\newcommand{\C}{\mathbb C}
\newcommand{\Q}{\mathbb Q}
\newcommand{\dist}{\operatorname{dist}}

\title{Countably Saturated Banach Spaces Have Countable Capture}
\author{Run \texttt{fa\_banach\_001}, lane 19}
\date{June 17, 2026}

\begin{document}
\maketitle

\section*{Source Problem}

The source paper is Jinghao Huang, Karimbergen Kudaybergenov, and Rui Liu,
\emph{The ball-covering property of von Neumann algebras and noncommutative
symmetric spaces}, arXiv:2311.14257.  After proving the commutative
\(C(K)\) criterion and the von Neumann algebra criterion, it asks for a
noncommutative \(C^*\)-algebra version:

\begin{center}
\includegraphics[width=0.92\textwidth]{figures/open_problem_crop.png}
\end{center}

\section*{Definitions}

Let \(X\) be a complex Banach space.  We use the following elementary
one-variable saturation principle.  Say that \(X\) is \emph{countably
norm-saturated} if every countable family of closed norm conditions in one
variable \(u\), with parameters from \(X\), that is finitely satisfiable in
the unit ball of \(X\), is realized in the unit ball of \(X\).

This is weaker than full countable saturation in continuous logic.  For
\(C^*\)-algebras, it is also weaker than the usual model-theoretic countable
saturation, and is covered by countable degree-one saturation for the specific
linear norm conditions used below.

\section*{Statement}

\begin{theorem}
\label{thm:saturated}
Let \(X\) be an infinite-dimensional countably norm-saturated complex Banach
space.  Then \(X\) has countable capture.  Consequently \(X\) fails BCP.
\end{theorem}

\begin{corollary}
\label{cor:cstar}
Every infinite-dimensional countably saturated \(C^*\)-algebra fails BCP.
The same conclusion holds under any saturation hypothesis that realizes
countable types consisting of the degree-one norm conditions in the proof.
\end{corollary}

\section*{Idea of the Proof}

For a finite set of centres \(F\subset X\), choose norming functionals
\(f_c\in J(c)\).  Since \(X\) is infinite dimensional, the finite
intersection \(\bigcap_{c\in F}\ker f_c\) contains a nonzero vector.  A unit
vector \(u\) in this intersection satisfies
\[
        \|c-\lambda u\|\ge |f_c(c-\lambda u)|=\|c\|
\]
for every scalar \(\lambda\).  Thus finite pieces of the countable-capture
requirements are always satisfiable.

For a countable centre set \(C\), write the requirements using only rational
complex scalars:
\[
        \|u\|=1,\qquad
        \|c-\lambda u\|\ge \|c\|
        \quad(c\in C,\ \lambda\in\Q+i\Q).
\]
Countable saturation realizes these simultaneously.  Continuity in
\(\lambda\) upgrades the inequalities from rational scalars to all complex
scalars, giving \(\dist(c,\C u)=\|c\|\) for every \(c\in C\).

\section*{Proof}

\begin{proof}[Proof of Theorem~\ref{thm:saturated}]
Let \(C=\{c_1,c_2,\ldots\}\subset X\) be countable.  Consider the following
countable type in the variable \(u\):
\[
        \|u\|=1,
\]
and, for every \(j\ge1\) and every \(\lambda\in\Q+i\Q\),
\[
        \|c_j-\lambda u\|\ge \|c_j\|.
\]
Equivalently, the latter can be written as the closed norm condition
\[
        \max\{0,\|c_j\|-\|c_j-\lambda u\|\}=0.
\]

We first prove finite satisfiability.  Fix finitely many centres
\(c_1,\ldots,c_N\).  Discard the zero centres.  For each remaining \(c_j\),
choose \(f_j\in X^*\) such that
\[
        \|f_j\|=1,\qquad f_j(c_j)=\|c_j\|.
\]
The space
\[
        \bigcap_{j=1}^N\ker f_j
\]
has finite codimension in the infinite-dimensional space \(X\), and is
therefore nonzero.  Choose \(u\) in this intersection with \(\|u\|=1\).  Then
for every listed centre and every \(\lambda\in\C\),
\[
        \|c_j-\lambda u\|
        \ge
        |f_j(c_j-\lambda u)|
        =
        \|c_j\|.
\]
The zero centres satisfy the same inequality trivially.  Hence every finite
part of the type is satisfiable.

By countable norm-saturation, there is \(u\in X\) realizing the whole type.
Then \(\|u\|=1\), and for every \(j\) and every rational complex
\(\lambda\),
\[
        \|c_j-\lambda u\|\ge\|c_j\|.
\]
For fixed \(j\), the map \(\lambda\mapsto \|c_j-\lambda u\|\) is continuous,
and \(\Q+i\Q\) is dense in \(\C\).  Therefore
\[
        \|c_j-\lambda u\|\ge\|c_j\|\qquad(\lambda\in\C).
\]
Taking the infimum over \(\lambda\) gives
\[
        \dist(c_j,\C u)\ge\|c_j\|.
\]
The reverse inequality follows by taking \(\lambda=0\).  Hence
\[
        \dist(c_j,\C u)=\|c_j\|\qquad(j\ge1).
\]
By the countable-capture characterization from the earlier lane-19 packet,
\(X\) has countable capture.  The countable-capture obstruction then implies
that \(X\) fails BCP.
\end{proof}

\begin{proof}[Proof of Corollary~\ref{cor:cstar}]
A countably saturated \(C^*\)-algebra is countably saturated as a metric
structure, so in particular it realizes the one-variable norm type used in
Theorem~\ref{thm:saturated}.  Apply the theorem to the underlying Banach
space.
\end{proof}

\section*{Scope for the C*-Classification}

This result adds a model-theoretic negative mechanism to the previous
topological, tensorial, product, von Neumann, group, and corona mechanisms.
It explains why ultrapower and reduced-power style \(C^*\)-algebras naturally
fall on the negative side: the finite centre obstruction is automatic in
infinite dimension, and saturation turns finite satisfiability into a single
global escaping direction.

The theorem does not settle arbitrary simple nonseparable \(C^*\)-algebras of
density continuum, because such algebras need not be countably saturated or
even countably degree-one saturated.

\section*{Novelty and Search Bounds}

The local run indexes were searched for saturation, countable capture,
ultrapowers, reduced powers, and countably degree-one saturation.  No prior
lane-19 packet recorded this general saturation-to-capture implication.
Web searches on June 17, 2026 found nearby saturation literature for
\(C^*\)-algebras and coronas, including Farah--Hart on countable saturation of
corona algebras, Eagle--Vignati on saturation and elementary equivalence of
\(C^*\)-algebras, and Voiculescu on countable degree-one saturation.

\section*{Human Review Focus}

Reviewers should check:
\begin{enumerate}
    \item the finite satisfiability argument via finitely many norming
          functionals;
    \item that the countable type uses only one-variable norm conditions;
    \item the passage from rational scalars to all complex scalars;
    \item the stated relationship between the custom saturation hypothesis and
          standard countable saturation/countable degree-one saturation.
\end{enumerate}

\begin{thebibliography}{9}

\bibitem{HKL}
Jinghao Huang, Karimbergen Kudaybergenov, and Rui Liu,
\emph{The ball-covering property of von Neumann algebras and noncommutative
symmetric spaces},
arXiv:2311.14257, 2023.

\bibitem{FH}
Ilijas Farah and Bradd Hart,
\emph{Countable saturation of corona algebras},
arXiv:1112.3898, 2011.

\bibitem{EV}
Christopher J. Eagle and Alessandro Vignati,
\emph{Saturation and elementary equivalence of \(C^*\)-algebras},
arXiv:1406.4875, 2014.

\bibitem{Voiculescu}
Dan-Virgil Voiculescu,
\emph{Countable Degree-1 Saturation of Certain C*-Algebras Which Are Coronas
of Banach Algebras},
arXiv:1310.4862, 2013.

\bibitem{lane19capture}
Run \texttt{fa\_banach\_001}, lane 19,
\emph{When the Countable-Capture BCP Obstruction Exists},
solution packet
\texttt{2311.14257\_countable\_capture\_characterization}, 2026.

\end{thebibliography}

\end{document}

